\subsection{Definitionen}

Dieser Abschnitt führt die grundlegenden Definitionen und Konzepte ein, die für das Verständnis von Counterfactual Regret Minimization notwendig sind.

\subsubsection{Wahrscheinlichkeiten und Werte}

Für ein Strategieprofil $\sigma$ definieren wir folgende Wahrscheinlichkeiten und Werte:

Die Wahrscheinlichkeit, dass eine History $h$ auftritt, wenn alle Spieler gemäß $\sigma$ handeln, bezeichnen wir mit $\pi^\sigma(h)$.
Diese Wahrscheinlichkeit kann in ein Produkt zerlegt werden:
\begin{equation}
\pi^\sigma(h) = \prod_{i \in N \cup \{c\}} \pi_i^\sigma(h),
\end{equation}
wobei $\pi_i^\sigma(h)$ die Wahrscheinlichkeit bezeichnet, dass Spieler $i$ in all seinen Entscheidungsknoten auf dem Pfad zu $h$ gemäß $\sigma$ die Aktionen gewählt hat, die zu $h$ führen.
Das Produkt der Wahrscheinlichkeiten aller anderen Spieler (inklusive Chance) bezeichnen wir mit $\pi_{-i}^\sigma(h)$.

Für ein Informationsset $I$ definieren wir die Erreichbarkeitswahrscheinlichkeit als
\begin{equation}
\pi^\sigma(I) = \sum_{h \in I} \pi^\sigma(h).
\end{equation}
Analog definieren wir $\pi_i^\sigma(I)$ und $\pi_{-i}^\sigma(I)$ als die Wahrscheinlichkeiten, dass Spieler $i$ beziehungsweise die anderen Spieler das Informationsset $I$ erreichen.

Der erwartete Payoff für Spieler $i$ unter Strategieprofil $\sigma$, auch \emph{Overall Value} genannt, berechnet sich als
\begin{equation}
u_i(\sigma) = \sum_{h \in Z} u_i(h) \pi^\sigma(h),
\end{equation}


\subsubsection{Nash Gleichgewicht}

Ein Nash Gleichgewicht für ein Zweispieler Nullsummenspiel ist ein Strategieprofil $\sigma = (\sigma_1, \sigma_2)$, welches folgende Ungleichungen erfüllt:
\begin{equation}
u_1(\sigma) \ge \max_{\sigma_1' \in \Sigma_1} u_1(\sigma_1', \sigma_2),
\qquad
u_2(\sigma) \ge \max_{\sigma_2' \in \Sigma_2} u_2(\sigma_1, \sigma_2')
\end{equation}
(\cite[Eq.~1]{zinkevich2007regret}).

Eine Approximation eines Nash Gleichgewichts, auch $\varepsilon$-Nash Gleichgewicht genannt, ist ein Strategieprofil $\sigma$, welches folgende Ungleichungen erfüllt:
\begin{equation}
u_1(\sigma) + \varepsilon \ge \max_{\sigma_1' \in \Sigma_1} u_1(\sigma_1', \sigma_2),
\qquad
u_2(\sigma) + \varepsilon \ge \max_{\sigma_2' \in \Sigma_2} u_2(\sigma_1, \sigma_2')
\end{equation}
(\cite[Eq.~2]{zinkevich2007regret}).
